% %%%%%%%%%%%%%%%%%%%%%%%%%%%%%%%%%%%%%%%%%%%%%%%%%%%%%%%%%%%%
% CHAPTER 5 - CONCLUSION
% %%%%%%%%%%%%%%%%%%%%%%%%%%%%%%%%%%%%%%%%%%%%%%%%%%%%%%%%%%%%
%
% DRAFT. Before submission:
%  - The findings summarised here depend on the review data in Chapter 4.
%    If your checks of the included studies change any judgement, update
%    Sections 5.1 and 5.2 accordingly.
%  - Add page numbers to the citations (SRH Guidelines 4.2.1).
%  - Rewrite the text in your own words.

\chapter{Conclusion}
\label{ch:conclusion}

This chapter concludes the thesis. Section~\ref{sec:summary} summarises the main findings, and Section~\ref{sec:answers} answers the four research questions. Section~\ref{sec:implications} discusses the theoretical and managerial implications, and Section~\ref{sec:futurelimits} reflects critically on the limitations of the thesis and outlines directions for future research.

% ------------------------------------------------------------
\section{Summary of Findings}
\label{sec:summary}

The thesis set out to evaluate how far published empirical evidence supports the proposition that AI capabilities contribute to a sustainable competitive advantage in e-commerce and platform-based businesses through dynamic capabilities. Chapter~\ref{ch:intro} argued that this question matters because AI tools are becoming widely accessible, so that any lasting advantage is more likely to depend on how firms organise around AI than on the technology itself. Building on the resource-based view and dynamic capabilities theory, Chapter~\ref{ch:litreview} developed a conceptual model in which AI capability relates to sustainable competitive advantage directly (H1), relates to dynamic capabilities (H2), and works through sensing, seizing and transforming individually (H3 to H5) and through dynamic capabilities as a whole (H6). Chapter~\ref{ch:method} described a structured literature review with explicit criteria, quality appraisal and decision rules to assess this model.

The review in Chapter~\ref{ch:analysis} identified 1,807 records and included 120 studies. Its findings can be condensed into four key messages. First, the literature is new and growing fast: 112 of the included studies were published in 2024 or later. Second, AI is consistently associated with competitive advantage, but mainly indirectly. In two of the appraised studies, the direct association of AI became small or non-significant once other capabilities, such as agility or AI-enabled sustainability capabilities, were included in the model \parencite{srahin2026, texeira2026}, although other studies, including the high-quality study of platform sellers, report direct and indirect associations side by side \parencite{guo2025}. Third, the specific mechanism proposed in this thesis has hardly been tested. Only 10 of the 120 studies address the dimensions of dynamic capabilities, only three of them quantitatively, and no study estimates separate indirect effects through sensing, seizing and transforming on competitive advantage. Fourth, the evidence from e-commerce and platform-based business is very thin: only four empirical studies concern this context, and the only one of high quality examines complementors that sell through cross-border platforms rather than platform owners or retailers \parencite{guo2025}.

Measured against the decision rules defined in Table~\ref{tab:rules}, no hypothesis meets the criteria for \textit{supported}, and all six are at most partly supported. The basis of this judgement differs greatly: H1 and H2 rest on several directly reported results from studies of mostly unknown quality, H3, H4 and H5 on a single indirectly relevant result each, and H6 on studies of related capabilities rather than of dynamic capabilities as a higher-order construct. This outcome does not mean that the proposed relationships are absent. It means that the published evidence is not yet strong, specific or context-relevant enough to confirm them.

The thesis makes three contributions. First, it provides a relationship-by-relationship assessment of a literature that has grown so quickly that its overall state is difficult to see: by separating the direct link, the link to dynamic capabilities and each mediating path, it shows that the support for the overall argument is much stronger than the support for its individual steps. Second, it makes the basis of each judgement explicit. The decision rules, the distinction between directly reported and indirectly relevant evidence, and the sensitivity checks allow readers to see exactly why a hypothesis is judged as it is and how the judgement would change if the evidence changed. Third, it develops a revised conceptual model that adds related capabilities and boundary conditions to the original model and marks the paths that remain untested, which turns the gaps in the evidence into a research agenda.

% ------------------------------------------------------------
\section{Answers to the Research Questions}
\label{sec:answers}

\textbf{RQ1: To what extent does published empirical evidence support a relationship between AI capabilities and sustainable competitive advantage?}

The evidence supports this relationship only in part. Nine independent studies report a positive association between AI capability and competitive advantage directly, although for six of them this classification is provisional, and 44 further studies report related evidence, for example with AI use instead of AI capability or with competitive advantage reached through other capabilities (Section~\ref{sec:h1}). The direct evidence, however, comes from studies whose quality is low or could not be assessed, and the one high-quality study coded as contradictory, in which AI capabilities were modelled mainly as a moderator, found no significant main effect \parencite{wang2026}. In two appraised studies and one study judged on its abstract, AI relates to competitive advantage mainly through mediating capabilities \parencite{srahin2026, texeira2026, chen2026}, but this is not the case in every study \parencite{guo2025}. In addition, only one study measured whether an advantage persists over time \parencite{dzreke2025}. The answer is also sensitive to the quality of the evidence: among the studies whose methods could be appraised, only one reports the relationship directly, and the AI construct that most studies measure is the extent of AI use rather than AI capability. The answer to RQ1 is therefore that published evidence supports an association between AI and competitive advantage, but not a direct relationship with an advantage that is demonstrably sustainable.

\textbf{RQ2: To what extent does published empirical evidence support a relationship between AI capabilities and dynamic capabilities?}

The evidence again supports this relationship in part. Three independent samples report directly that AI capability is positively related to dynamic capabilities \parencite{almheiri2025, fano2025, neirafernandez2026}, and 42 studies relate AI to capabilities that resemble one dimension of dynamic capabilities, such as agility, innovation capability, absorptive capacity or market sensing (Section~\ref{sec:h2}). The direction of the relationship is, however, unresolved. Some studies treat dynamic capabilities as an antecedent of AI adoption \parencite{deng2025}, and others treat AI as a moderator of the effect of dynamic capabilities \parencite{quoc2026}. Because almost all studies are cross-sectional, they cannot decide between these specifications. Moreover, all directly reported evidence comes from studies that could only be judged on their abstracts, and from settings such as the public sector, in which competition plays little role. The answer to RQ2 is that AI and dynamic capabilities are consistently associated, but the evidence does not show that AI capability builds dynamic capabilities rather than the reverse.

\textbf{RQ3: To what extent does published empirical evidence indicate that sensing, seizing and transforming capabilities each mediate the relationship between AI capabilities and sustainable competitive advantage?}

The published evidence hardly addresses this question. For sensing, a single study of e-commerce firms reports an indirect effect of generative AI adoption on marketing performance through market sensing \parencite{noor2026}, which makes H3 partly supported on the basis of one indirectly relevant study. For seizing and transforming, a single study of tourism businesses reports indirect effects of AI orientation on firm performance through absorptive capacity and opportunity recognition, which its authors interpret as seizing and transforming mechanisms \parencite{li2026}; under the same rule, H4 and H5 are therefore partly supported on the basis of one indirectly relevant study each. No study estimates an indirect effect through sensing, seizing or transforming, measured as dimensions of dynamic capabilities, on competitive advantage. Of the ten studies that address the three dimensions explicitly, four are conceptual, one is a Delphi study and two are qualitative or mixed-methods, so that theorising about the dimensions clearly runs ahead of testing. The answer to RQ3 is therefore that the dimension-specific mediation proposed by dynamic capabilities theory \parencite{teece2007} has not yet been tested in a way that allows a conclusion.

\textbf{RQ4: To what extent does published empirical evidence indicate that dynamic capabilities, as a higher-order construct, mediate this relationship?}

The evidence indicates that organisational capabilities mediate the relationship, but not specifically dynamic capabilities as a higher-order construct. Fifty-one studies provide indirectly relevant evidence, most of them estimating indirect effects through related capabilities such as innovation capability, agility, international marketing capabilities or organisational learning \parencite{malik2021, srahin2026, guo2025, tu2026}. No study estimates the indirect effect of AI capability on competitive advantage through dynamic capabilities as defined in this thesis, none specifies dynamic capabilities as a higher-order construct, one study finds no mediation by dynamic capabilities \parencite{kokash2026}, and some studies merge AI capability and dynamic capabilities into a single construct \parencite{ayoub2026, shao2026}. In addition, only 13 studies estimate an indirect effect directly, whereas 38 infer mediation from separate paths. H6 is therefore partly supported only in the sense defined by the decision rules, through related capabilities. The answer to RQ4 is that the mediating logic of the conceptual model is broadly consistent with the evidence, but that the mediator is better described as a family of capabilities strengthened by AI than as dynamic capabilities in the narrow sense.

% ------------------------------------------------------------
\section{Implications}
\label{sec:implications}

The findings of a review have implications for theory and for practice, but these implications have to match the strength of the evidence. Because none of the hypotheses is fully supported, the implications of this thesis are stated conditionally: they describe what follows if the observed patterns hold, and they identify where theory and practice should be cautious. The following subsections discuss the theoretical implications first and the managerial implications second.

\subsection{Theoretical Implications}
\label{sec:theoryimp}

The findings have three implications for theory. First, they support the resource-based argument that AI on its own is unlikely to be a source of sustained advantage. Because AI is valuable but widely accessible, it rarely meets the criteria of rareness and inimitability \parencite{barney1991}. The finding of several studies that its direct association weakens once firm-specific capabilities are considered is consistent with this view and with the argument that AI must be situated in the organisation to create an advantage \parencite{kemp2024}. It also mirrors earlier evidence on big data analytics, whose effect on competitive performance was fully mediated by dynamic capabilities \parencite{mikalef2020}.

Second, the findings qualify the role of dynamic capabilities in this mechanism. The literature does not test sensing, seizing and transforming as distinct mediators; it tests a wide range of more specific capabilities. This is consistent with the view that dynamic capabilities take the form of identifiable processes, such as product development or strategic decision making \parencite{eisenhardt2000}, but it leaves open whether these processes add up to the higher-order capability described by \textcite{teece2007}. Future theorising should therefore specify how particular capabilities, such as agility or absorptive capacity, relate to the three dimensions, instead of using dynamic capabilities as a general label.

Third, the findings suggest that AI may amplify existing capabilities as much as it creates new ones. Several studies found that AI strengthened the effect of other capabilities or orientations rather than acting directly \parencite{quoc2026, wang2026}, and that its effect depended on innovation orientation, leadership or environmental turbulence \parencite{texeira2026, sofiyan2026, kokash2026}. This interpretation fits the idea that AI changes the sources of competitive advantage through complementation with existing human and organisational capabilities \parencite{krakowski2023}. The revised model in Figure~\ref{fig:revised} therefore adds related capabilities and boundary conditions to the original model.

A fourth implication concerns the clarity of the constructs. The review found that AI capability, AI use and AI integration are often treated as interchangeable, that some studies define AI capability as a dynamic capability in its own right \parencite{shao2026, ayoub2026}, and that dynamic capabilities are frequently invoked as a theoretical lens without being measured. Under these conditions, a positive association between \enquote{AI} and \enquote{capabilities} can mean very different things. Theoretical progress therefore depends less on new hypotheses than on agreed definitions and measures: AI capability as a firm-specific bundle of resources \parencite{mikalef2021ai}, dynamic capabilities as the higher-order capacity to sense, seize and transform \parencite{teece2007}, and competitive advantage as a position relative to competitors that persists over time \parencite{barney1991}. The debate about situated AI \parencite{kemp2024, moser2024} adds that the embedding of AI in organisational routines may be the decisive feature, and this feature is not captured by current survey measures at all.

\subsection{Managerial Implications}
\label{sec:managerialimp}

For managers in e-commerce and platform-based businesses, the main implication is that investment in AI technology is unlikely to produce a lasting advantage on its own. The evidence indicates that AI creates value through capabilities that surround it, such as the ability to act quickly on new information, to innovate and to learn \parencite{srahin2026, malik2021, tu2026}. Managers should therefore assess AI projects not only by their technical performance but also by whether they strengthen these organisational capabilities, and they should invest in the processes, skills and structures that allow AI-generated insights to be turned into decisions and changes.

A second implication concerns the conditions under which AI pays off. The indirect effect of AI increased with a sustainable innovation orientation in one study of restaurants \parencite{texeira2026} and depended on transformational leadership in another \parencite{sofiyan2026}. Leadership attention and a strategic orientation towards innovation are therefore part of an AI strategy rather than separate issues. For firms that sell through platforms, the evidence on sellers using cross-border platforms suggests that AI is most effective when it is combined with strategic agility and marketing capabilities \parencite{guo2025}.

The implications differ between types of firms in digital markets. Platform owners control the data generated on their platforms and can therefore use AI to strengthen data network effects, but research suggests that this advantage depends on legitimacy, data stewardship and user-centric design \parencite{gregory2021}. AI capability can even reduce the value users perceive if it raises concerns about privacy \parencite{yu2026}, and opaque, poorly governed AI can lead to poor decisions and competitive disadvantage \parencite{rana2022}. For platform owners, AI governance is therefore part of the competitive strategy, not merely a compliance task. E-commerce retailers, by contrast, compete with rivals who can buy similar AI tools, so their advantage depends on how well AI is combined with their own customer knowledge, marketing capabilities and capacity to change \parencite{noor2026, bunea2026}. Complementors that sell through platforms rely partly on the platform's AI; for them, the evidence suggests that combining AI with their own strategic agility and marketing capabilities is what creates an advantage over other sellers \parencite{guo2025}.

Finally, managers should treat claims about the competitive benefits of AI with caution. Most of the evidence comes from cross-sectional surveys of managers' perceptions, often outside e-commerce, and the effects of AI may take time to appear in performance, because AI requires complementary intangible investments \parencite{brynjolfsson2021}. Firms that monitor the effects of their AI investments over several years, and compare them with competitors, will be better placed to judge whether an advantage is lasting.

% ------------------------------------------------------------
\section{Limitations and Future Research}
\label{sec:futurelimits}

The conclusions of this thesis must be read in the light of its limitations, which were set out in Section~\ref{sec:limitations}. The most important are that the review relied on a single database, that the citation searching could not be completed, that the quality of 75 of the 95 empirical studies could not be assessed, mostly because only their abstracts were accessible, that nine included studies could be judged only on their titles, and that screening and extraction were supported by a language model. The thesis also inherits the limitations of the primary studies, most of which are cross-sectional surveys of single respondents and therefore vulnerable to common method bias and reverse causality \parencite{podsakoff2024, hill2021}. The judgements in Chapter~\ref{ch:analysis} are therefore statements about the current state of the published evidence, not about the true effect of AI. Repeating the review with additional databases, the completed citation search and full access to the included studies would strengthen these judgements.

A critical reflection on the research process adds three further points. First, the decision to exclude studies that relate AI only to firm performance was made after a scoping search and before screening. It kept the review focused on the dependent variable and the mechanism of the thesis, but it also removed a large body of evidence on the performance effects of AI, including well-known studies \parencite{mikalef2021ai, chen2022}, so a broader review might have reached somewhat different conclusions on H1. Second, the use of a language model made it possible to screen almost 1,800 records twice and to require a verbatim quotation from the accessible source text for every extracted result, which would have been difficult for a single researcher. At the same time, it introduces a new kind of risk, because the model may misread studies in systematic ways that are harder to detect than individual human errors; the manual checks described in Section~\ref{sec:criteria} are therefore essential. Third, the decision rules proved strict. Because the quality of most studies could not be assessed, no hypothesis could reach the judgement \textit{supported}. This strictness was intended, because it prevents the volume of weak evidence from being mistaken for strength, but it also means that the judgements describe the limits of the evidence as much as the relationships themselves.

These limitations, together with the gaps identified in Chapter~\ref{ch:analysis}, point to five directions for future research. First, studies should measure sensing, seizing and transforming as distinct dimensions and estimate the indirect effect through each of them, so that H3 to H5 can be tested directly. Higher-order models of dynamic capabilities should be specified and validated with established procedures \parencite{sarstedt2019}. Second, research needs longitudinal or time-lagged designs, ideally combining survey data with archival measures of performance over several years, to establish the direction of the relationships and the persistence of an advantage. Third, the constructs need to be kept apart: AI capability should be measured separately from dynamic capabilities and from AI adoption, so that their relationships can be examined rather than assumed.

Fourth, the context of e-commerce and platform-based business deserves dedicated study. Mechanisms that are specific to platforms, such as data network effects \parencite{gregory2021} and the integrative capabilities through which platform leaders orchestrate their ecosystems \parencite{helfat2018}, have not yet been examined empirically, and the differences between platform owners, retailers and complementors remain unexplored \parencite{cennamo2021}. Fifth, the boundary conditions found in this review, such as innovation orientation, leadership and environmental turbulence, should be built into future models from the start. Qualitative case studies of platform firms could help to identify how these conditions shape the conversion of AI capability into a lasting advantage.

Two further directions follow from the critical discussion in Chapter~\ref{ch:analysis}. Research should give more attention to the negative and null effects of AI. Only one of the 70 extracted results on AI and competitive advantage was negative, which is unlikely to reflect the full range of experiences with AI in practice, given the risks of opaque and poorly governed AI \parencite{rana2022}. Registered reports, studies that explicitly test for failure and the publication of null results would give a more balanced picture. Finally, the rapid spread of generative AI raises the question of whether the findings of studies conducted before its diffusion still apply. Generative AI lowers the cost of many capabilities that were previously firm-specific, such as content creation and customer communication, and may therefore make AI-based advantages even harder to sustain. Future research should examine whether generative AI capabilities follow the same indirect logic as earlier forms of AI or whether they erode the advantages that firms have built.

In conclusion, the thesis finds that AI capability is associated with competitive advantage mainly through the organisational capabilities that AI strengthens, but that the evidence is not yet sufficient to confirm that dynamic capabilities, and in particular sensing, seizing and transforming, carry this effect in e-commerce and platform-based businesses. Answering this question will require research that is more specific in its constructs, longer in its time horizon and closer to the context of digital platforms.
